\subsection[The residual pairs (2,b), b >= 6]{The residual pairs $(2,b)$, $b\ge6$}\label{loc:S5}
Let $P=h_2f_b$, $m=b-2\ge4$, $R_0=K^2G_m(K)g$.
For $0\le x\le\theta$ the finite residue formulas give
\[
 |G_m(i\sqrt x)|^2\le(1+x^3)^2/(1+x)\le1-2x/5.
\]
For the second inequality it suffices that
$2x^2+x^5+2x/5\le3/5$; its left side is at most
$25/72+1/32+1/6=157/288$.
Consequently
\begin{gather*}
 \|KG_m(K)g\|^2\le \delta_{2}-2\delta_{3}/5,\quad
 \|R_0\|^2\le \delta_{3}-2\delta_{4}/5,\\
 \E[KG_m(K)g]\ge-d+\delta_{2}-\delta_{3},\quad
 \langle g,R_0\rangle\ge-\delta_{2}+\delta_{3}-\delta_{4},\quad
 \E[pR_0]=-u_{b+1}\le \delta_{5}\le\theta^4d.
\end{gather*}
Here $u_j=\E K^j\one$. The lower bounds retain one further negative term in each finite
alternating sum. The last identity is exact telescoping; its first possible
positive value occurs at index $b+1=10$.

Set $a=19/32$, $\epsilon=ad$, $\tau=\theta^4=28561/1048576$,
$m_0=1-ad$ and $R=R_0-\epsilon g$.
Retaining $f_0,f_1$ in convolution gives
$h_2+f_b\le2m_0+R$, so $P\le(m_0+R/2)^2$ and
$\E[pR]\le L=\tau d-ad^2$.
The exact row and variance budgets are
\begin{align*}
 \|R\|_\infty^2&\le \delta_{2}-2\delta_{3}/5+\epsilon^2+2\epsilon(d-\delta_{2}+\delta_{3})
                       \le U=871d/2560+20539d^2/16384,\\
 \E R^2&\le \delta_{3}-2\delta_{4}/5+\epsilon^2d+2\epsilon(\delta_{2}-\delta_{3}+\delta_{4})\le dW,\\
 W&=11323/81920+191919d/524288+361d^2/1024.
\end{align*}
For the two spectral substitutions the derivatives of
$(1-2\epsilon)x+(2\epsilon-2/5)x^2$ and
$x^2-2x^3/5+2\epsilon(x-x^2+x^3)$ are nonnegative on $[0,\theta]$.
Thus $U,W$ increase, and $V_p/d\le W+\sqrt{UW}$.
The cap is $P\le(19/18)^2$: the square of $2ad+1/9$ minus $U$ is
\[
 2565d^2/16384-1759d/23040+1/81>0,
\]
whose discriminant is $-1009919/530841600$.
In particular $(19/18)^6<7/5$.

The common binomial estimates have main terms
\[
 H=(1-ad)^{11}[1-(a-6\tau)d-6ad^2],\qquad
 F=(1-ad)^5[1-(a-3\tau)d-3ad^2].
\]
They are convex on $[0,1/3]$. For $H$, put $b=a-6\tau$, $c=6a$;
the bracket in $H''/(1-ad)^9$ is
$110a^2+22ab-2c+12a(4c-11ab)d-156a^2cd^2$.
Its linear term is positive and the remainder is at least
$a(110a-104a^2-12)+22ab>0$.
For $F$, put $b=a-3\tau$, $c=3a$; the corresponding bracket is
$20a^2+10ab-2c+6a(4c-5ab)d-42a^2cd^2$,
bounded below by $2a(10a-7a^2-3)+10ab>0$.
These signs use $1/2<a<3/5$ and $d^2\le1/9$.
For $d\le1/32$ the exponential criterion gives
$D-3U>d[1/8-13d/4]\ge3d/128$, with
$D\ge(2a-\tau)d+ad^2$.

The remaining sixth-moment intervals use $H+B_{12}V_p$ and the following
complete data; all bounds are strict at positive endpoints.
\begin{center}\small
\begin{tabular}{@{}ccccc@{}}\toprule
$I$&$\bar m$&$M$&$V_p/d$&$B_{12}$\\\midrule
$[1/32,1/16]$&$53/54$&$1/6$&$1/4$&$19$\\
$[1/16,3/32]$&$26/27$&$5/24$&$11/40$&$17$\\
$[3/32,1/8]$&$17/18$&$1/4$&$3/10$&$31/2$\\
$[1/8,5/32]$&$25/27$&$29/100$&$47/140$&$27/2$\\
$[5/32,11/64]$&$49/54$&$5/16$&$9/25$&$12$\\\bottomrule
\end{tabular}
\end{center}
The endpoint upper bounds are respectively
\begin{gather*}
 (377/384,185/192),\quad(1841/1920,1265/1280),\quad
 (631/640,157/160),\\
 (433/448,5375/5376),\quad(29/30,799/800).
\end{gather*}
For the cubic intervals use $F+B_6V_p$:
\begin{center}\small
\begin{tabular}{@{}ccccc@{}}\toprule
$I$&$\bar m$&$M$&$V_p/d$&$B_6$\\\midrule
$[11/64,3/16]$&$44/49$&$1/3$&$19/50$&$63/20$\\
$[3/16,7/32]$&$8/9$&$11/30$&$5/12$&$31/10$\\
$[7/32,1/4]$&$67/77$&$13/32$&$23/50$&$3$\\
$[1/4,2/7]$&$109/128$&$9/20$&$41/80$&$113/40$\\\bottomrule
\end{tabular}
\end{center}
The respective endpoint upper bounds are
$(228203/320000,11091/16000)$,
$(91/128,81/200+217/768)$,
$(1131/1600,139/200)$ and $(9113/12800,7993/11200)$.
All are less than $5/7$, so the pointwise cap converts the cubic to the
sixth moment. Each table uses the increasing $U,W$, the full finite sum
$B_r$, and convex interpolation of the displayed main polynomial.

Finally, for $d\ge2/7$, the row bound $|r|\le pq/2$ gives
$ph_2\le pq(1+p/2)\le37/24$; for $g\ge0$ use $g-3g^2\le1/12$,
and for $g\le0$ use $3/2$ directly.
Positivity gives $f_b\le2\mu_{b-2}p$, hence $P\le37\mu_{b-2}/12$.
For $b\ge7$ the proved mean bound $\mu_{b-2}\le1-19d/8$ gives $111/112$.
For $b=6$ retain the exact $\mu_4\le1-83d/32$, giving $1073/1344$.
Both are less than one. This covers all $d$ and proves $\E[pP^6]\le1$.
